\documentclass{article}
\usepackage[utf8]{inputenc}
\usepackage{graphicx}
\usepackage{amsmath}
\usepackage{amssymb}
\usepackage[margin = 0.5in]{geometry}
\usepackage{collectbox}

\linespread{1.8}

\makeatletter
\newcommand{\mybox}{%
    \collectbox{%
        \setlength{\fboxsep}{1pt}%
        \fbox{\BOXCONTENT}%
    }%
}
\makeatother


\title{MAT 527:  Functions of a Complex Variable I \\
	   \textbf{Submission I}}
\author{Nolan R. H. Gagnon}
\begin{document}
\maketitle
\large
\noindent\mybox{\textbf{Theorem}} If $\left\lbrace U_{\lambda} \right\rbrace_{\lambda\in\Lambda}$ is a family of open sets, then $\bigcup_{\lambda\in\Lambda} U_{\lambda}$ is open.

\vspace{3mm}

\noindent\textbf{Proof:}  Let $\left\lbrace U_{\lambda} \right\rbrace_{\lambda\in\Lambda}$ be a family of open sets, where $\Lambda$ is some index set.  Let $z\in \bigcup_{\lambda\in\Lambda} U_{\lambda}$.  Then for some $\mu\in\Lambda$, $z$ is an element of $U_{\mu}$.  Now, since $U_{\mu}$ is open, there exists $\Delta >0$ such that $B(z;\Delta)\subseteq U_{\mu}$.  But clearly $U_{\mu}\subseteq \bigcup_{\lambda\in\Lambda} U_{\lambda}.$  Therefore, $B(z;\Delta)\subseteq \bigcup_{\lambda\in\Lambda} U_{\lambda}.$  Therefore, $\bigcup_{\lambda\in\Lambda} U_{\lambda}$ is open.
\hfill$\blacksquare$
\end{document}
